\section{Точные импульсы: периферия RMT}\label{les:rmt}
\lessonintro
  {формировать пакеты уровней аппаратно и выбирать точность времени}
  {уроки 17 и 21}
  {пакет из двух импульсов с проверенной длительностью}
  {логический анализатор или вторая плата}

\subsection{Когда программная задержка мешает}
Для адресного светодиода или ИК-пульта важна длительность каждого уровня. Цикл с задержками
зависит от прерываний и планировщика. RMT получает массив уровней и длительностей, затем
формирует или измеряет фронты аппаратно. Периодический PWM удобнее получать LEDC,
а переменный пакет --- RMT.

Один символ содержит две части: \code{level0/duration0} и \code{level1/duration1}.
При 1 МГц один такт равен 1 мкс, при 10 МГц --- 0,1 мкс. Более высокая точность уменьшает
максимальный интервал при той же разрядности поля. У ESP32-S3 TX и RX имеют отдельные ресурсы;
библиотека NeoPixel уже может занимать RMT, поэтому учитывайте соседние драйверы.

\subsection{Пакет, который можно проверить}
Соедините GPIO4 с входом анализатора, земли объедините. Без анализатора можно соединить GPIO4
с GPIO5 опыта PCNT из урока 31. Счётчик увидит два положительных фронта на пакет,
но не проверит их длительность. На выход подаётся логика 3,3 В.

\begin{figure}[H]
\centering
\begin{circuitikz}[scale=0.95,transform shape]
  \node[blk,fill=blkcpu,minimum width=2.8cm,minimum height=4.2cm] (m) at (0,0) {};
  \node[blk,fill=blkper,minimum width=2.8cm,minimum height=2.6cm] (d) at (8,0.5) {};
  \node[font=\sffamily\small\bfseries] at (0,2.6) {ESP32-S3};
  \node[font=\sffamily\small\bfseries,align=center] at (8,2.2) {Логический\\[-2pt]анализатор};
  \foreach \y/\n in {1/{\pin{4}},0/{\pin{5}},-1.5/GND}
    \node[font=\sffamily\scriptsize,anchor=east] at (1.4,\y) {\n};
  \foreach \y/\n in {1/CH0,-0.4/GND}
    \node[font=\sffamily\scriptsize,anchor=west] at (6.6,\y) {\n};
  \draw[thick,accent] (1.4,1) -- (6.6,1);
  \node[font=\sffamily\scriptsize,text=accent,fill=white,inner sep=1pt] at (4.9,1) {выход RMT};
  \draw[thick,black!70] (1.4,-1.5) -- (5.8,-1.5) -- (5.8,-0.4) -- (6.6,-0.4);
  % вариант без анализатора
  \draw[thick,dashed,warn] (1.4,0) -- (2.6,0) -- (2.6,1);
  \fill[accent] (2.6,1) circle (1.5pt);
  \node[font=\sffamily\scriptsize,text=warn,anchor=west,align=left] at (2.9,0.3) {без анализатора:\\[-2pt] перемычка на \pin{5}};
\end{circuitikz}
\caption{Проверка выхода RMT: \pin{4} подключается к входу логического анализатора, земли
объединяются. Пунктиром показана замена анализатору --- перемычка на вход PCNT из урока 31}
\end{figure}

\codecap{Две пары уровней: 100/200 и 300/400 мкс}
\codeinput{l30_rmt_pulses}

\begin{figure}[H]\centering
\begin{tikzpicture}[x=0.011cm,y=0.7cm]
\draw[->] (0,0)--(1100,0) node[right,font=\small]{мкс};
\draw[thick,esp] (0,0)--(0,1)--(100,1)--(100,0)--(300,0)--(300,1)--(600,1)--(600,0)--(1000,0);
\foreach \x/\t in {0/0,100/100,300/300,600/600,1000/1000}
\node[below,font=\scriptsize] at (\x,0) {\t};
\end{tikzpicture}
\caption{Один пакет RMT, общая длительность 1000 мкс}
\end{figure}

\subsection{ИК: несущая, огибающая и команда}
ИК-пульт обычно передаёт внутри отметки колебания несущей, например 38 кГц.
Демодулирующий приёмник убирает несущую и выдаёт огибающую, нередко активным LOW.
Передатчику нужны согласованная несущая, драйвер светодиода и ограничение тока.
Наш пример создаёт логический пакет, а не готовую команду телевизору.

Для декодирования сначала выделяют старт, затем длительности битов и паузы; конкретные
числа берут из описания выбранного протокола. Начните с записи собственного пульта:
сравните одиночное нажатие и удержание, найдите повторяющиеся участки. Ошибка приёма
не всегда означает плохую библиотеку: источником может быть световая помеха или питание.

\subsection{Асинхронность и память}
\code{rmtWrite} ждёт завершения с ограниченным тайм-аутом. При асинхронном вызове массив
должен жить до окончания передачи. Локальный массив завершившейся функции для этого не подходит;
также нельзя менять буфер, который ещё читает периферия.
При RX конец пакета определяют паузой: слишком короткий порог разрежет сообщение,
слишком длинный задержит его обработку. Учитывайте ёмкость буфера и число символов.

\begin{trouble}
Нет сигнала: неверный GPIO или занят канал. Пакет обрезан: не хватает буфера либо тайм-аута.
ИК-приёмник молчит: нет нужной несущей или перепутаны питание и выход.
Проверяйте сначала форму сигнала, затем смысл команды.
\end{trouble}
\begin{summary}
RMT задаёт аппаратные длительности, частота счётчика определяет цену такта.
Несущая, огибающая и содержание команды --- разные уровни описания.
\end{summary}

\begin{selfcheck}
\begin{enumerate}
\item Сколько длятся 250 тактов при 1 МГц?
\item Почему PCNT недостаточно для проверки пакета?
\item Как долго должен жить буфер асинхронной передачи?
\end{enumerate}
\end{selfcheck}
\needspace{20\baselineskip}
\begin{tasks}
\begin{enumerate}
\item \lv{1} Пересчитайте четыре длительности примера в микросекунды.
\item \lv{1} Запишите пакет и сравните с рисунком.
\item \lv{2} Добавьте третий импульс 500/500 мкс и измерьте общую длительность.
\item \lv{2} Сравните точность и предельный интервал при 1 и 10 МГц.
\item \lv{3} Передавайте асинхронно, продолжая читать кнопку; обеспечьте время жизни буфера.
\item \lv{4} Запишите огибающую собственного ИК-пульта, выделите старт, биты и повтор удержания.
\end{enumerate}
\end{tasks}

\paragraph{Документация для продолжения}\href{https://docs.espressif.com/projects/arduino-esp32/en/latest/api/rmt.html}{Официальное руководство по теме}.
